%% pc-triangle-feu — le triangle du feu : combustible, comburant, énergie d'activation.
%% Les trois composantes doivent être présentes SIMULTANÉMENT : supprimer un seul côté
%% (rôle de l'extincteur, variante grisée de droite) suffit à éteindre le feu.
\documentclass[tikz,border=4pt]{standalone}
\usepackage{liaisons}
\usetikzlibrary{shapes.geometric, positioning}
\begin{document}
\begin{tikzpicture}[x=1cm,y=1cm,
  cote/.style={line width=2pt, line cap=round, line join=round, draw=solideD},
  etiq/.style={font=\small\bfseries, text=solideD, inner sep=3pt},
  ex/.style={font=\scriptsize, text=black!65, inner sep=1pt}]

%% flamme stylisée centrée en (0,0), hauteur ~1 cm : \flamme{couleur ext}{couleur int}
\newcommand\flammechemin{(0,0.60)
    .. controls (0.10,0.30) and (0.34,0.24) .. (0.28,-0.06)
    .. controls (0.24,-0.32) and (-0.24,-0.32) .. (-0.28,-0.06)
    .. controls (-0.34,0.22) and (-0.07,0.26) .. (0,0.60)}
\newcommand\flamme[2]{%
  \draw[draw=#1, line width=1pt, fill=#1!18, line join=round] \flammechemin;
  \draw[draw=#2, line width=0.8pt, fill=#2!35, line join=round]
    (0,0.26) .. controls (0.06,0.10) and (0.17,0.06) .. (0.14,-0.08)
             .. controls (0.11,-0.22) and (-0.11,-0.22) .. (-0.14,-0.08)
             .. controls (-0.17,0.08) and (-0.04,0.10) .. (0,0.26);}

%% ================= triangle du feu (complet) ===============================
\coordinate (Bg) at (-1.55,0);      % sommet bas gauche
\coordinate (Bd) at ( 1.55,0);      % sommet bas droit
\coordinate (S)  at ( 0,2.68);      % sommet haut
\draw[cote] (Bg) -- (S) -- (Bd) -- cycle;

\node[etiq, rotate=60,  anchor=south, align=center] at (-0.775,1.34)
  {Combustible\\[-1pt]\textcolor{black!65}{\scriptsize\mdseries essence, gaz, bois}};
\node[etiq, rotate=-60, anchor=south, align=center] at ( 0.775,1.34)
  {Comburant\\[-1pt]\textcolor{black!65}{\scriptsize\mdseries dioxygène $\mathrm{O_2}$ de l'air}};
\node[etiq, anchor=north, align=center] at (0,-0.10)
  {Énergie d'activation\\[-1pt]\textcolor{black!65}{\scriptsize\mdseries étincelle, flamme}};

\begin{scope}[shift={(0,1.10)}]  \flamme{solideD}{solideA} \end{scope}

\node[font=\scriptsize, text=black!70, align=center, anchor=north] at (0,-1.15)
  {les trois composantes doivent être\\présentes \textbf{simultanément}};

%% ================= variante : action d'un extincteur =======================
\begin{scope}[shift={(4.05,0.35)}, scale=0.62]
  \coordinate (bg) at (-1.55,0);
  \coordinate (bd) at ( 1.55,0);
  \coordinate (s)  at ( 0,2.68);
  \draw[cote, draw=black!30] (bg) -- (s) -- (bd) -- cycle;
  %% un côté supprimé : croix rouge sur le combustible
  \draw[line width=2.4pt, line cap=round, draw=solideA]
    (-1.12,1.06) -- (-0.44,1.62)  (-1.12,1.62) -- (-0.44,1.06);
  \begin{scope}[shift={(0,1.10)}]
    \draw[draw=black!40, line width=1pt, dashed, dash pattern=on 2pt off 1.6pt, line join=round]
      \flammechemin;
  \end{scope}
\end{scope}
\node[font=\scriptsize, text=solideA, align=center, anchor=north] at (4.05,0.20)
  {un côté supprimé\\(extincteur) :\\le feu s'éteint};
\end{tikzpicture}
\end{document}
